% !TEX program = xelatex
% Lernplatt - by Can Siebert
% Kurs 71 (Dig 42) - Tag 8 - Lehrskript
% Enhancement, Predictive Analytics, Real-time Monitoring, DWH/ETL, BI-Dashboards
%
% Self-contained xelatex source. Kein biblatex, keine externen Bilder.

\documentclass[11pt,a4paper,oneside]{article}

% --- Encoding & Sprache --------------------------------------------------
\usepackage{polyglossia}
\setdefaultlanguage{german}

% --- Geometrie -----------------------------------------------------------
\usepackage[a4paper,top=2cm,bottom=2cm,outer=2.5cm,inner=2.5cm,bindingoffset=1.5cm]{geometry}

% --- Fonts ---------------------------------------------------------------
\usepackage{fontspec}
\defaultfontfeatures{Ligatures=TeX}

\IfFontExistsTF{Charter}{%
  \setmainfont{Charter}[%
    Numbers=OldStyle,%
    UprightFeatures   = {SmallCapsFeatures={Letters=SmallCaps}}%
  ]%
}{%
  \IfFontExistsTF{Bitstream Charter}{%
    \setmainfont{Bitstream Charter}%
  }{%
    \setmainfont{TeX Gyre Schola}%
  }%
}

\IfFontExistsTF{Source Sans 3}{%
  \setsansfont{Source Sans 3}%
}{%
  \IfFontExistsTF{Source Sans Pro}{%
    \setsansfont{Source Sans Pro}%
  }{%
    \setsansfont{TeX Gyre Heros}%
  }%
}

\newcommand{\caveatfontfallback}{\itshape}
\IfFontExistsTF{Caveat}{%
  \newfontfamily\caveatfont{Caveat}[Scale=1.15]%
}{%
  \newcommand\caveatfont{\caveatfontfallback}%
}

\setmonofont{Latin Modern Mono}

% --- Mikro-Typografie ----------------------------------------------------
\usepackage{microtype}
\linespread{1.15}

% --- Farben & Boxen ------------------------------------------------------
\usepackage{xcolor}
\definecolor{lpInk}{RGB}{20,28,38}
\definecolor{kurs71}{HTML}{9D4A1A}
\definecolor{lpAccent}{HTML}{9D4A1A}
\definecolor{lpAccentLight}{RGB}{248,232,218}
\definecolor{lpAmber}{RGB}{222,138,30}
\definecolor{lpAmberLight}{RGB}{255,243,217}
\definecolor{lpAmberBorder}{RGB}{196,118,18}
\definecolor{lpGray}{RGB}{96,104,114}
\definecolor{lpGrayLight}{RGB}{238,240,243}

\usepackage[most]{tcolorbox}
\tcbuselibrary{breakable,skins}

\newtcolorbox{eselsbox}[1][Eselsbrücke]{%
  enhanced, breakable,
  colback=lpAmberLight,
  colframe=lpAmberBorder,
  boxrule=0.6pt,
  arc=2pt,
  borderline={0.4pt}{0pt}{lpAmberBorder, dashed},
  left=10pt,right=10pt,top=8pt,bottom=8pt,
  fonttitle=\sffamily\bfseries\color{lpAmberBorder},
  coltitle=lpAmberBorder,
  title={#1},
  attach boxed title to top left={xshift=8pt,yshift=-8pt},
  boxed title style={size=small, colback=lpAmberLight, colframe=lpAmberBorder, boxrule=0.4pt, arc=1pt},
  top=14pt
}

\newtcolorbox{merkbox}[1][Merksatz]{%
  enhanced, breakable,
  colback=lpAccentLight,
  colframe=lpAccent,
  boxrule=0.6pt,
  arc=2pt,
  left=10pt,right=10pt,top=8pt,bottom=8pt,
  fonttitle=\sffamily\bfseries\color{white},
  coltitle=white,
  title={#1},
  attach boxed title to top left={xshift=8pt,yshift=-8pt},
  boxed title style={size=small, colback=lpAccent, colframe=lpAccent, boxrule=0.4pt, arc=1pt},
  top=14pt
}

% --- Listen --------------------------------------------------------------
\usepackage{enumitem}
\setlist{itemsep=2pt, topsep=4pt, parsep=0pt}
\setlist[itemize,1]{label={\textbullet},leftmargin=1.6em}
\setlist[itemize,2]{label={\textendash},leftmargin=1.4em}
\setlist[enumerate,1]{label=\arabic*., leftmargin=1.8em}

% --- Tabellen ------------------------------------------------------------
\usepackage{tabularx}
\usepackage{array}
\usepackage{booktabs}
\usepackage{longtable}

% --- Kopf-/Fusszeilen ----------------------------------------------------
\usepackage{fancyhdr}
\usepackage{lastpage}
\pagestyle{fancy}
\fancyhf{}
\renewcommand{\headrulewidth}{0.3pt}
\renewcommand{\footrulewidth}{0pt}
\fancyhead[L]{\footnotesize\sffamily\color{lpGray}Lernplatt \textperiodcentered{} Kurs 71 \textperiodcentered{} Tag 8}
\fancyhead[R]{\footnotesize\sffamily\color{lpGray}Predictive \textperiodcentered{} DWH \textperiodcentered{} Dashboards}
\fancyfoot[C]{\footnotesize\sffamily\color{lpGray}\thepage{} / \pageref{LastPage}}

\fancypagestyle{plain}{%
  \fancyhf{}%
  \renewcommand{\headrulewidth}{0pt}%
  \fancyfoot[C]{\footnotesize\sffamily\color{lpGray}\thepage{} / \pageref{LastPage}}%
}

% --- Section-Styling -----------------------------------------------------
\usepackage[explicit]{titlesec}

\titleformat{\section}[block]
  {\sffamily\Large\bfseries\color{lpAccent}}
  {\thesection.}{0.6em}{#1}
  [{\vspace{2pt}\color{lpAccent}\titlerule[0.6pt]}]

\titleformat{\subsection}[block]
  {\sffamily\large\bfseries\color{lpInk}}
  {\thesubsection}{0.6em}{#1}

\titleformat{\subsubsection}[block]
  {\sffamily\normalsize\bfseries\color{lpInk}}
  {}{0pt}{#1}

\titlespacing*{\section}{0pt}{14pt}{8pt}
\titlespacing*{\subsection}{0pt}{10pt}{4pt}
\titlespacing*{\subsubsection}{0pt}{8pt}{2pt}

% --- Hyperlinks ----------------------------------------------------------
\usepackage[hidelinks,unicode]{hyperref}

% --- TOC -----------------------------------------------------------------
\renewcommand{\contentsname}{\sffamily\Large\bfseries\color{lpAccent}Inhalt}

% --- Helfer --------------------------------------------------------------
\newcommand{\akzent}[1]{{\caveatfont\color{lpAmber}#1}}
\newcommand{\fachbegriff}[1]{\textbf{#1}}
\newcommand{\quelle}[1]{{\footnotesize\color{lpGray}(#1)}}

% =========================================================================
\begin{document}

% --- COVER ---------------------------------------------------------------
\begin{titlepage}
  \pagestyle{empty}
  \vspace*{1.5cm}
  \noindent{\sffamily\footnotesize\color{lpGray}LERNPLATT \textperiodcentered{} BY CAN SIEBERT}\\[2pt]
  \noindent{\color{lpAccent}\rule{\linewidth}{1.2pt}}\\[4pt]
  \noindent{\sffamily\footnotesize\color{lpGray}MODUL 3 \textperiodcentered{} TAG 8 VON 2 \textperiodcentered{} KURS 71 (Dig 42) \textperiodcentered{} 9.\,Juni 2026}

  \vspace{3cm}
  {\sffamily\fontsize{30}{34}\selectfont\bfseries\color{lpInk}Predictive\\\&\\Dashboards\par}

  \vspace{0.7cm}
  {\sffamily\large\color{lpGray}Vom Ist verstehen zum Zukunft steuern --\\Enhancement, Vorhersagen und Real-time Monitoring auf DWH/BI-Basis.\par}

  \vspace{1.5cm}
  {\caveatfont\LARGE\color{lpAmber}Lehrskript -- Tag 8 von 16}\par

  \vfill

  \noindent\begin{minipage}{0.55\linewidth}
    {\sffamily\footnotesize\color{lpGray}Lernpfad:}\\
    {\sffamily\small\color{lpInk}Wissen \textrightarrow{} Anwendung \textrightarrow{} Management}\\[10pt]
    {\sffamily\footnotesize\color{lpGray}Bearbeitungsdauer:}\\
    {\sffamily\small\color{lpInk}1 Kurstag (8 Unterrichtseinheiten)}
  \end{minipage}\hfill
  \begin{minipage}{0.4\linewidth}
    \raggedleft
    {\sffamily\footnotesize\color{lpGray}Dozent:}\\
    {\sffamily\small\color{lpInk}Can Siebert}\\[10pt]
    {\sffamily\footnotesize\color{lpGray}Format:}\\
    {\sffamily\small\color{lpInk}Theorie \textperiodcentered{} Fallstudie \textperiodcentered{} Coaching}
  \end{minipage}

  \vspace{0.5cm}
  \noindent{\color{lpAccent}\rule{\linewidth}{0.6pt}}
\end{titlepage}

% --- LERNZIELE -----------------------------------------------------------
\section*{Lernziele dieses Tages}
\addcontentsline{toc}{section}{Lernziele dieses Tages}
\thispagestyle{plain}

\noindent Gestern haben wir mit Discovery und Conformance den Ist-Prozess gegen das Soll-Modell gespiegelt. Heute machen wir den nächsten Schritt: vom \emph{Verstehen} zum \emph{Vorhersagen} und \emph{Steuern in Echtzeit}. Drei Bausteine machen das möglich -- Process Enhancement (Modelle mit Daten anreichern), Predictive Process Analytics (Risiken vor Eintritt erkennen) und ein solides Daten-Fundament aus DWH, ETL und BI-Dashboards.

\subsection*{L1 -- Wissen (Reproduktion und Einordnung)}
\begin{itemize}
  \item Process Enhancement als Anreicherungs-Disziplin einordnen (Zeiten, Ressourcen, Kosten, Varianten).
  \item ML-Typen im Prozesskontext unterscheiden: Klassifikation, Regression, Clustering, Anomalieerkennung.
  \item DWH/ETL-Grundlagen kennen: Fakten und Dimensionen, Staging, Mapping, Quality Checks.
  \item Dashboard-Prinzipien benennen: Executive-KPIs, Drilldown, Alerts, ``next best action''.
\end{itemize}

\subsection*{L2 -- Anwendung (Transfer auf konkrete Situationen)}
\begin{itemize}
  \item Einen Predictive-Use-Case-Canvas (Zielentscheidung, Prognose, Handlung, Datenminimum, Risiko, Erfolg) ausfüllen.
  \item Eine Baseline-Heuristik formulieren, die ohne ML sofort wirken kann.
  \item Eine ETL-Pipeline in sechs Schritten skizzieren und fünf Data-Quality-Checks definieren.
  \item Ein Dashboard mit Executive-, Prozess- und Risiko-KPI sowie einer Alert-Kachel entwerfen.
\end{itemize}

\subsection*{L3 -- Management (Entscheidung und Verantwortung)}
\begin{itemize}
  \item Predictive Analytics als \emph{Produkt} verantworten: Product Owner, Release-Zyklus, Modell-Monitoring.
  \item Alerts und Playbooks koppeln: jeder Alarm hat Owner, Handlung, Eskalation und Dokumentation.
  \item Anti-Gaming-Mechanismen bei Executive-KPIs etablieren.
  \item Entscheidungen unter Unsicherheit treffen: welche Use Cases zuerst, welche Datenlücken bewusst akzeptiert.
\end{itemize}

\begin{merkbox}[Roter Faden]
Eine Vorhersage ohne Handlung ist Reporting. Ein Alert ohne Owner ist Lärm. Ein Modell ohne Monitoring altert unbemerkt. Senior-Praxis verbindet alle drei: \emph{Prognose $\rightarrow$ Owner $\rightarrow$ Playbook $\rightarrow$ Feedback in das Modell.}
\end{merkbox}

\clearpage

% --- 1. EINSTIEG ---------------------------------------------------------
\section{Einstieg: Vom Rückblick zur Zukunftssteuerung}

Die meisten BPM-Programme haben in den letzten zehn Jahren gelernt, Vergangenheit zu \emph{messen}: KPIs aus Vormonats-Daten, Reports aus Quartals-Abschlüssen, Audit-Findings aus Stichproben. Was darin fehlt, ist die Fähigkeit zur \emph{Zukunftssteuerung}: \emph{Welche Cases werden in den nächsten Stunden eskalieren? Welcher Engpass bildet sich gerade? Welche Anomalie sollte heute mein Eingreifen auslösen?} \quelle{Davenport, 2018; Maggi et al., 2014}

\subsection{Drei Reifegrade der Prozesssteuerung}

\begin{enumerate}
  \item \fachbegriff{Reporting-Reifegrad.} Vergangenheit beschreiben (Monatsbericht, Quartals-KPI). Hilft bei Diagnose, nicht bei Steuerung.
  \item \fachbegriff{Real-time-Reifegrad.} Gegenwart beobachten (Live-Dashboard, KPI in Near-Real-Time). Schnellere Diagnose, aber noch reaktiv.
  \item \fachbegriff{Predictive-Reifegrad.} Zukunft vorhersagen (Welche Cases werden in den nächsten 4\,h SLA reißen?). Erlaubt \emph{präventives} Eingreifen.
\end{enumerate}

\subsection{Die drei Bausteine des heutigen Tages}

\begin{itemize}
  \item \fachbegriff{Process Enhancement.} Modelle werden mit Daten angereichert -- Zeiten, Ressourcen, Kosten, Varianten -- und werden so vom Bild zum steuerbaren Modell.
  \item \fachbegriff{Predictive Process Analytics + Real-time Monitoring.} ML-Modelle erzeugen Vorhersagen (SLA-Breach, Rework, Abbruch), die mit Alerts und Playbooks verknüpft werden.
  \item \fachbegriff{DWH/ETL und BI-Dashboards.} Das Datenfundament, ohne das Predictive und Monitoring beliebig werden.
\end{itemize}

\begin{eselsbox}[Eselsbrücke: \akzent{E-P-D}]
Drei Bausteine, in Reihenfolge des Wertbeitrags:\\[2pt]
\textbf{E}nhancement -- Modelle mit Daten anreichern.\\
\textbf{P}redictive -- Zukunft anhand der Daten antizipieren.\\
\textbf{D}ashboards -- Steuerung und Handlung sichtbar machen.\\[2pt]
\emph{Wer ohne E startet, baut ML auf Sand. Wer ohne D arbeitet, hat Modelle, die niemand sieht.}
\end{eselsbox}

% --- 2. PROCESS ENHANCEMENT ---------------------------------------------
\section{Process Enhancement -- vom Bild zum steuerbaren Modell}

\fachbegriff{Process Enhancement} ist die dritte Disziplin neben Discovery und Conformance. Während Discovery aus Logs ein Modell erzeugt und Conformance Ist und Soll vergleicht, \emph{erweitert} Enhancement vorhandene Modelle um Daten: Zeiten, Kosten, Ressourcen, Varianten, Kontext. Das Modell wird dadurch nicht mehr nur \emph{strukturell}, sondern \emph{operationell} verwertbar. \quelle{van der Aalst, 2016}

\subsection{Drei Hauptdimensionen der Anreicherung}

\begin{itemize}
  \item \fachbegriff{Zeit.} Bearbeitungszeit (BT), Wartezeit (WT), Durchlaufzeit (DLZ), p50/p95-Quartile pro Schritt.
  \item \fachbegriff{Kosten und Ressourcen.} Welche Ressource hat den Schritt gemacht, was kostet er, wer ist Engpass?
  \item \fachbegriff{Varianten und Kontext.} Welche Pfad-Varianten gibt es, korreliert mit welchem Kontext (Standort, Produkt, Kunden\-segment, Tageszeit)?
\end{itemize}

\subsection{Beispiel -- ein Modell wird operationell}

Ein einfaches BPMN-Modell ``Reklamation bearbeiten'' wird durch Enhancement zu:

\begin{itemize}
  \item Activity ``Klassifizieren'': BT-Median 12\,min, WT-Median 35\,min, p95 BT 45\,min, p95 WT 4\,h.
  \item Engpass: Klassifizieren tagsüber Personalengpass, abends $<$5 Cases/h.
  \item Variante: ``Klassifizieren $\rightarrow$ Rückfrage $\rightarrow$ Klassifizieren'' tritt in 18\,\% der Cases auf, korreliert mit Produktkategorie X.
  \item Kosten: durchschnittlich 8\,\text{\euro}\ pro Case, davon 65\,\% Personalkosten Klassifizieren.
\end{itemize}

Aus einem ``Bild'' wird so eine Steuerungsbasis: \emph{wo lohnt sich Personal-Aufstockung? Wo Re-Design? Wo Automatisierung?}

\subsection{Voraussetzungen für Enhancement}

\begin{itemize}
  \item Event Log mit Pflicht-Spalten (C-A-T, siehe Tag 7).
  \item Idealerweise zusätzlich Ressource, Kosten, Kontext-Attribute.
  \item Saubere Datenqualität (Mapping, Case-ID-Strategie, Dedup).
  \item Soll-Modell als Vergleichsbasis.
\end{itemize}

\begin{merkbox}[Enhancement schließt den BPM-Lifecycle]
Discovery + Conformance + Enhancement bilden zusammen die analytische Schicht eines BPM-Programms. Ohne Enhancement bleiben Discovery-Modelle hübsche Bilder; ohne Conformance bleiben Enhancement-Daten reine Statistik; ohne Discovery fehlt der Anschluss an die Realität.
\end{merkbox}

% --- 3. ML IM PROZESSKONTEXT --------------------------------------------
\section{Machine Learning in der Prozessoptimierung}

Predictive Process Monitoring -- also ML-basierte Vorhersagen auf Prozessdaten -- ist die jüngste Schicht des Process-Mining-Spektrums. Die meisten Use Cases lassen sich in vier ML-Typen einordnen. \quelle{Maggi et al., 2014; Shmueli, 2010}

\subsection{Vier ML-Typen für Prozesse}

\begin{enumerate}
  \item \fachbegriff{Klassifikation.} Vorhersage einer binären oder kategorialen Zielgröße: \emph{Wird Case eskalieren?} (ja/nein), \emph{Welcher Pfad wird gewählt?} (A/B/C).
  \item \fachbegriff{Regression.} Vorhersage einer kontinuierlichen Größe: \emph{Wie lange wird Case noch dauern? Wie hoch werden Kosten?}.
  \item \fachbegriff{Clustering.} Gruppierung ähnlicher Cases ohne vorgegebene Labels: \emph{Welche Falltypen gibt es überhaupt?}.
  \item \fachbegriff{Anomalieerkennung.} Identifikation ungewöhnlicher Pfade, Fehler oder potenzieller Manipulation: \emph{Was unterscheidet diesen Case stark vom Rest?}.
\end{enumerate}

\subsection{Typische Predictive Use Cases}

\begin{itemize}
  \item \fachbegriff{SLA-Risiko.} ``Wahrscheinlichkeit, dass dieser Case die SLA verletzt, in den nächsten 4\,h.''
  \item \fachbegriff{Rework-Risiko.} ``Wahrscheinlichkeit, dass dieser Case in eine Korrekturschleife geht.''
  \item \fachbegriff{Abbruch-/Storno-Risiko.} ``Wahrscheinlichkeit, dass dieser Case nicht erfolgreich abgeschlossen wird.''
  \item \fachbegriff{Engpassbildung.} ``Welche Queue wird in den nächsten 2\,h überlastet sein?''
\end{itemize}

\subsection{Grenzen und Risiken}

\begin{itemize}
  \item \fachbegriff{Datenqualität.} Schlechte Daten erzeugen schöne, aber falsche Modelle.
  \item \fachbegriff{Konzeptdrift.} Modelle, die auf alten Daten trainiert wurden, verlieren Genauigkeit, wenn sich die Realität ändert.
  \item \fachbegriff{Erklärbarkeit.} ``Black-Box''-Modelle (Deep Learning, manche Ensembles) sind schwer zu erklären -- problematisch im regulierten Umfeld.
  \item \fachbegriff{Gaming/Fehlanreize.} Wenn Mitarbeitende wissen, wie das Modell entscheidet, können sie es manipulieren (Datenfeld so füllen, dass Alert nicht auslöst).
  \item \fachbegriff{Bias/Fairness.} Wenn das Modell historische Diskriminierung lernt, verlagert es sie in die Vorhersagen.
\end{itemize}

\subsection{Baseline vor ML}

Eine pragmatische Heuristik: \emph{Bevor man ML einsetzt, prüft man, ob eine einfache Baseline-Regel 60--80\,\% des Nutzens bringt.} Beispiele:

\begin{itemize}
  \item ``Wenn Ticket-Alter $> 70\,\%$ der SLA-Zeit, dann Alert.''
  \item ``Wenn Anzahl Reassignments $> 2$, dann Eskalation.''
  \item ``Wenn Service-Kategorie X UND Schicht Y, dann Vorabkennzeichen ``kritisch''.''
\end{itemize}

Baseline-Regeln haben drei Vorteile: \emph{auditierbar, erklärbar, sofort wirksam.} Sie sind keine ``ML-Schande'', sondern oft die produktivste erste Stufe.

\begin{merkbox}[Baseline schlägt schlechtes ML]
Eine erklärbare Baseline-Regel mit 70\,\%-Genauigkeit ist oft wertvoller als ein ML-Modell mit 80\,\%-Genauigkeit, das niemand versteht. Wer mit Baseline startet und gezielt nach Lücken sucht, baut ML auf Substanz. Wer mit ML startet, baut oft auf einen Datenstand, der dafür nicht reif ist.
\end{merkbox}

\subsection{Predictive Use Case Canvas}

Ein praktikabler Canvas zur Strukturierung eines Predictive Use Case:

\begin{enumerate}
  \item \textbf{Zielentscheidung.} Was soll konkret entschieden werden? (z.\,B.\ ``Soll dieser Case priorisiert werden?'')
  \item \textbf{Prognose.} Was wird vorhergesagt? (z.\,B.\ ``SLA-Breach-Wahrscheinlichkeit in 4\,h'')
  \item \textbf{Handlung.} Was passiert bei hohem Risiko? (z.\,B.\ ``Duty Manager Notification + Triage'')
  \item \textbf{Datenminimum.} Welche 5--10 Felder werden zwingend benötigt?
  \item \textbf{Risiko/Compliance.} Was darf nicht passieren? (Falsch-Negative, Diskriminierung, Datenschutz)
  \item \textbf{Erfolgskriterium.} Wie messen wir den Nutzen? (z.\,B.\ ``SLA-Hit-Rate steigt um 10\,\%'')
\end{enumerate}

\subsubsection{Canvas für ``SLA-Risiko früh erkennen''}

\begin{itemize}
  \item Ziel: Priorisierung bei knappen Ressourcen.
  \item Prognose: P(SLA-Breach in den nächsten 4\,h).
  \item Handlung: Risk $>0,8$ $\rightarrow$ Duty Manager Notification; Risk $0,5$--$0,8$ $\rightarrow$ Markierung im Dashboard.
  \item Datenminimum: Priorität, Ticket-Alter, Anzahl Reassignments, Service-Kategorie, Schichtzeit.
  \item Risiko: False Positives erzeugen Alarmmüdigkeit; Datenschutz bei Personenbezug.
  \item Erfolg: SLA-Hit-Rate $\geq$95\,\% bei stabilem oder reduzierten Personaleinsatz.
\end{itemize}

\subsubsection{Canvas für ``Rework früh erkennen''}

\begin{itemize}
  \item Ziel: Frühe Korrektur statt späte Wiedereröffnung.
  \item Prognose: P(Rework innerhalb von 7 Tagen).
  \item Handlung: Risk $>0,7$ $\rightarrow$ Vier-Augen-Prüfung vor Abschluss; Risk $0,4$--$0,7$ $\rightarrow$ Checkliste an Bearbeitenden.
  \item Datenminimum: Falltyp, Kategorie, Anzahl Korrekturschleifen, Bearbeitungsdauer, Kunden\-segment.
  \item Risiko: Über-Eingreifen verlangsamt den Prozess.
  \item Erfolg: Rework-Rate halbiert sich bei stabilem Bearbeitungs-Tempo.
\end{itemize}

% --- 4. REAL-TIME MONITORING --------------------------------------------
\section{Real-time Monitoring und Alert Playbooks}

Predictive ohne Real-time-Anbindung bleibt akademisch. Erst die Verknüpfung von Vorhersage, Anzeige, Alert und Handlungsanweisung macht aus Predictive Analytics ein Steuerungsinstrument. \quelle{Few, 2013; Davenport, 2018}

\subsection{Vier Bausteine eines Real-time-Programms}

\begin{itemize}
  \item \fachbegriff{Event Stream.} Quelle der Daten in Near-Real-Time (Kafka, Cloud-Bus, Change Data Capture).
  \item \fachbegriff{Inference Engine.} ML-Modell oder Baseline-Regel, die Vorhersagen erzeugt.
  \item \fachbegriff{Alert-System.} Schwellenwert-basierte Auslöser, mit Eskalationsstufen.
  \item \fachbegriff{Playbook.} Schriftliche Handlungsanweisung: wer macht was, in welcher Zeit, mit welcher Dokumentation.
\end{itemize}

\subsection{Alert Design}

\subsubsection{Schwellen festlegen}

\begin{itemize}
  \item Drei Stufen: \emph{warn} (gelb), \emph{alert} (orange), \emph{critical} (rot).
  \item Schwellen anhand historischer Daten kalibrieren, nicht raten.
  \item Bei Unsicherheit: konservativ starten (eher zu wenig Alerts), mit definiertem Re-Tuning.
\end{itemize}

\subsubsection{Alert-Ermüdung verhindern}

\begin{itemize}
  \item Pro Empfänger \emph{maximal} Alerts pro Tag/Schicht festlegen.
  \item Korrelation: zehn Alerts ``derselbe Service down'' nicht als zehn Alerts schicken.
  \item Snooze-/Acknowledge-Mechanismus, damit Bearbeitete nicht erneut alarmieren.
  \item Regelmäßig (z.\,B.\ monatlich) Alert-Wirksamkeit reviewen: welche Alerts führten zu Handlung?
\end{itemize}

\subsection{Playbook-Struktur}

Ein gutes Playbook für einen kritischen Alert hat fünf Komponenten:

\begin{enumerate}
  \item \textbf{Owner.} Wer ist verantwortlich (Rolle, Kontaktweg)?
  \item \textbf{Handlung.} Was ist die erste Handlung in den ersten Minuten?
  \item \textbf{Eskalation.} Wann und an wen wird eskaliert, falls Handlung nicht greift?
  \item \textbf{Dokumentation.} Wo wird die Bearbeitung dokumentiert (Ticket, Audit-Log)?
  \item \textbf{Post-Review.} Wer schaut bei wiederkehrenden Mustern auf Ursachen-Analyse?
\end{enumerate}

\subsection{Beispiel-Playbook ``SLA-Breach-Risiko hoch''}

\begin{itemize}
  \item \textbf{Trigger:} Risk $>0,8$ für Case mit Priorität High.
  \item \textbf{Owner:} Duty Manager Schicht aktiv.
  \item \textbf{Handlung:} ``Fast Triage'' in 15\,min: Case lesen, mit Bearbeitenden abstimmen, ggf.\ Ressourcen umschichten.
  \item \textbf{Eskalation:} Falls Risk weiterhin $>0,8$ nach 30\,min, an Head of Operations.
  \item \textbf{Dokumentation:} Eintrag im Ticket-System, Risk-Snapshot, Maßnahme.
  \item \textbf{Post-Review:} Bei wiederholtem Auftreten (3 Cases pro Woche): Service-Kategorie analysieren, Verbesserung priorisieren.
\end{itemize}

\begin{eselsbox}[Eselsbrücke: \akzent{O-H-E-D-P}]
Fünf Pflichtfelder jedes Alert-Playbooks:\\[2pt]
\textbf{O}wner -- wer trägt?\\
\textbf{H}andlung -- was sofort?\\
\textbf{E}skalation -- wann an wen?\\
\textbf{D}okumentation -- wo wird festgehalten?\\
\textbf{P}ost-Review -- wer lernt?\\[2pt]
\emph{Wenn ein Buchstabe fehlt, ist es kein Playbook, sondern eine Hoffnung.}
\end{eselsbox}

% --- 5. DWH UND ETL -----------------------------------------------------
\section{DWH und ETL als Daten-Fundament}

Predictive Analytics, Conformance und Dashboards stehen alle auf demselben Fundament: einem \fachbegriff{Data Warehouse (DWH)}, in dem Prozessdaten aus verschiedenen Quellsystemen konsolidiert, gemappt und qualitätsgesichert werden. Der Prozess der Befüllung heißt \fachbegriff{ETL} (Extract -- Transform -- Load) oder in moderner Variante \fachbegriff{ELT} (Extract -- Load -- Transform), bei der die Transformation im DWH selbst geschieht. \quelle{Kimball \& Ross, 2013; Inmon, 2005; Russom, 2011}

\subsection{Typische Quellsysteme}

\begin{itemize}
  \item \fachbegriff{ERP} (Aufträge, Rechnungen, Stammdaten).
  \item \fachbegriff{CRM} (Kunden, Kommunikation, Interaktionen).
  \item \fachbegriff{Ticket-/Workflow-Systeme} (Cases, Bearbeitungsschritte, Zuweisungen).
  \item \fachbegriff{Logs aus Workflow-Engines} (BPMS-Audit-Trail).
  \item \fachbegriff{Externe Daten} (Marktdaten, Wetter, Lieferanten-Updates).
\end{itemize}

\subsection{Sechs ETL-Schritte}

\begin{enumerate}
  \item \fachbegriff{Extract.} Aus Quellsystemen via API, DB-Replikation oder Datei-Export entnehmen.
  \item \fachbegriff{Staging.} Rohdaten unverändert in einem Staging-Bereich ablegen, bevor sie transformiert werden -- Sicherheitsnetz für Re-Runs.
  \item \fachbegriff{Mapping.} Aktivitäten harmonisieren (``Approve''/``Genehmigen''/``OK''), Lookup-Tabellen pflegen.
  \item \fachbegriff{Join.} Datenströme zu Events mit Case ID verknüpfen.
  \item \fachbegriff{Quality Checks.} Vor dem Laden ins DWH systematisch prüfen.
  \item \fachbegriff{Load DWH.} Faktentabelle ``Events'' plus Dimensionen befüllen.
\end{enumerate}

\subsection{Datenmodell -- Fakten und Dimensionen}

In der \fachbegriff{Dimensionalen Modellierung} (Kimball) gibt es zwei Tabellenklassen:

\begin{itemize}
  \item \fachbegriff{Faktentabelle Events} -- jede Zeile ein Event, mit Maßzahlen (Bearbeitungs\-zeit, Kosten) und Fremdschlüsseln zu den Dimensionen.
  \item \fachbegriff{Dimensionen} -- beschreibende Attribute: Aktivität, Zeit, Kunde, Produkt, Standort, Ressource. Jede Dimension hat eine ID und beschreibende Felder.
\end{itemize}

Vorteil dieses Modells: einfaches Star-Schema, schnelle Aggregationen, gute Eignung für BI-Tools.

\subsection{Fünf Data-Quality-Checks}

\begin{enumerate}
  \item \fachbegriff{Fehlende Case ID.} Events ohne Case ID lassen sich nicht zuordnen $\rightarrow$ in Quarantäne, manuelle Klärung.
  \item \fachbegriff{Zeitstempel rückwärts.} Event vor seinem Vorgänger laut Logik $\rightarrow$ Hinweis, Datenquelle prüfen.
  \item \fachbegriff{Unbekannte Aktivität.} Activity-Name nicht in Mapping-Tabelle $\rightarrow$ Mapping ergänzen oder Datenquelle prüfen.
  \item \fachbegriff{Dubletten.} Identische Event-Signaturen $\rightarrow$ Dedup-Regel anwenden, Zähler im Quality-Report.
  \item \fachbegriff{Plausibilität.} Bearbeitungs\-zeit $> N$ Standardabweichungen vom Mittel $\rightarrow$ Outlier-Markierung, optional Ausschluss.
\end{enumerate}

\begin{merkbox}[ETL ist 60\,\% des Wertes]
In jedem Mining-/Predictive-Programm liegt der Großteil der Wert\-schöpfung in der ETL-Schicht: Mapping, Case-ID-Strategie, Quality Checks. Tools für Discovery, Conformance und ML kommen und gehen -- ein sauberes, gepflegtes DWH ist die Investition, die jahrelang trägt.
\end{merkbox}

% --- 6. BI DASHBOARDS ----------------------------------------------------
\section{BI Dashboards -- vom Datenstand zur Steuerung}

Ein Dashboard ist kein Selbstzweck. Es ist ein \emph{Steuerungswerkzeug} für definierte Adressaten und Entscheidungen. Drei Prinzipien helfen, gute Dashboards von dekorativen zu unterscheiden. \quelle{Few, 2013}

\subsection{Drei Dashboard-Prinzipien}

\begin{enumerate}
  \item \fachbegriff{Executive KPIs (wenige, klar).} Auf der obersten Ebene maximal 6--8 Kacheln. Jede mit Ziel/Schwelle, klarer Beschriftung, Trend.
  \item \fachbegriff{Drilldown (Ursachen).} Ein Klick auf eine Kachel öffnet die nächsten Detailebenen -- Segmentierung nach Standort, Produkt, Zeit, Schicht.
  \item \fachbegriff{Alarme (Handlung).} Eine Kachel zeigt aktive Alerts mit Owner und nächstem Schritt -- damit das Dashboard nicht nur informiert, sondern Handlung ermöglicht.
\end{enumerate}

\subsection{Was ein gutes Dashboard zeigt}

\begin{itemize}
  \item \fachbegriff{Trend.} Wo bewegen wir uns hin (Linie, Sparkline)?
  \item \fachbegriff{Ziel/Schwellen.} Soll-Wert, Warn- und Kritisch-Linien.
  \item \fachbegriff{Segmentierung.} Mindestens eine Aufschlüsselung (Standort, Service, Kunden\-segment).
  \item \fachbegriff{Next best action.} Was ist die empfohlene Reaktion?
\end{itemize}

\subsection{Beispiel-Dashboard ``SLA-Frühwarnsystem''}

Sechs Kacheln:

\begin{enumerate}
  \item \textbf{Executive 1 -- SLA-Hit-Rate Tag} (Zielwert 95\,\%, Trend 7 Tage).
  \item \textbf{Executive 2 -- Anzahl aktive Cases mit Risk $>0,5$}.
  \item \textbf{Prozess 1 -- Top-10 Risikotickets} (Risk, Service, Owner).
  \item \textbf{Prozess 2 -- Heatmap Service $\times$ Queue} (Risk-Konzentration).
  \item \textbf{Risk/Compliance -- Alarm Log} (Alerts der letzten 24\,h, Status, Owner).
  \item \textbf{Alert-Kachel -- aktuelle kritische Alerts} (Risk $>0,8$, Eskalations\-status).
\end{enumerate}

\subsection{Anti-Gaming -- welche KPI \emph{nicht} sofort}

Eine bewusste Entscheidung in jedem Dashboard-Design ist die \emph{Auslassung}. Beispiele für problematische KPIs:

\begin{itemize}
  \item \fachbegriff{Cases pro Mitarbeitenden pro Stunde.} Risiko: ``Schnell-und-Schlampig'' wird belohnt. Gegenmaßnahme: Kopplung an Rework-Rate, oder gar nicht einführen.
  \item \fachbegriff{Erstanruf-Lösungsquote.} Risiko: Tickets werden vorzeitig geschlossen. Gegenmaßnahme: Kopplung an Reopen-Rate.
  \item \fachbegriff{Bearbeitungs\-zeit-Reduktion.} Risiko: ``Quick-Close'' ohne Lösung. Gegenmaßnahme: Kombination mit Kunden-NPS.
\end{itemize}

% --- 7. FALLSTUDIE -------------------------------------------------------
\section{Fallstudie: FinServe IT Operations -- SLA-Frühwarnsystem}

\emph{FinServe IT Operations} ist stark SLA-getrieben -- jedes Service-Level-Agreement gegenüber den Geschäftsbereichen ist vertraglich gesichert. Problem: SLA-Verletzungen werden zu spät erkannt, Eskalationen passieren hektisch, Management will Real-time-Transparenz.

\subsection{Ausgangslage und Restriktionen}

\begin{itemize}
  \item Zeit: 8 Wochen.
  \item Budget: 120\,000\,\text{\euro}.
  \item Daten: Ticketing, Monitoring, E-Mail.
  \item Labels (SLA breach) sind unvollständig.
  \item Zielkonflikt: Transparenz vs.\ Druck auf Teams.
\end{itemize}

\subsection{Schritt 1 -- Predictive Use Case definieren}

\textbf{Vorhersage:} ``Wahrscheinlichkeit SLA breach in den nächsten 4\,h.''
\textbf{Entscheidung:} Priorisierung der Tickets bei Bearbeitungs-Engpass.
\textbf{Handlung:} Bei Risk $>0,8$ Duty Manager benachrichtigen + Triage 15\,min.

\subsection{Schritt 2 -- Features (mindestens 8)}

\begin{enumerate}
  \item Priorität (P1--P4).
  \item Ticket-Alter (Minuten).
  \item Queue (Service-Gruppe).
  \item Anzahl Reassignments.
  \item Service betroffen (kategorial).
  \item Kategorie (Incident-Typ).
  \item Open Incidents per Service (Last-Indikator).
  \item Time since last update.
  \item Shift (Tag/Nacht/Wochenende).
  \item Customer tier (Gold/Silver/Bronze).
\end{enumerate}

\subsection{Schritt 3 -- Label und Umgang mit Lücken}

\textbf{Label:} SLA breach = Ticket nicht gelöst vor SLA-Zeit.

\textbf{Lücken:} Historische Labels nur teilweise gepflegt.

\textbf{Pragmatische Lösung:} Baseline-Heuristik ``Ticket Age $>$ 70\,\% SLA'' als Proxy-Label fürs erste Training; parallel manuelle Label-Pflege, in Release 2 ablösen.

\subsection{Schritt 4 -- ETL-Pipeline}

\begin{enumerate}
  \item \textbf{Extract.} Ticketing API + Monitoring Stream + E-Mail-Headers.
  \item \textbf{Staging.} Rohdaten 30 Tage rollend.
  \item \textbf{Mapping.} Service/Kategorie auf Standardliste, harmonisierte Activity-Namen.
  \item \textbf{Join.} Events mit Monitoring-Signalen über Service-ID.
  \item \textbf{Quality Checks.} Case ID, Timestamps, Dubletten, Plausibilität, Mapping-Coverage.
  \item \textbf{Load DWH.} Faktentabelle Events + Dimensionen (Ticket-Facts, Service, Customer, Time).
\end{enumerate}

\subsection{Schritt 5 -- Dashboard und Playbook}

\textbf{Dashboard-Kacheln:}
\begin{enumerate}
  \item SLA-Risk Heatmap (Service $\times$ Queue).
  \item Top-10 Risikotickets.
  \item Trend Breach-Rate 7\,d.
  \item Durchschnittliche Resolution Time.
  \item Anzahl Reassignments.
  \item Alert Log (mit Owner, Status).
\end{enumerate}

\textbf{Playbook:}
\begin{itemize}
  \item Trigger: Risk $>0,8$.
  \item Owner: Duty Manager.
  \item Handlung: ``Fast Triage'' in 15\,min, Ressourcen-Umschichtung.
  \item Eskalation: bei anhaltendem Risk an Head of Operations.
  \item Dokumentation: im Ticket, Risk-Snapshot, Maßnahmen-Log.
  \item Post-Review: monatlicher Service-Excellence-Termin.
\end{itemize}

\subsection{Schritt 6 -- MVP-Entscheidung}

\textbf{Im MVP:}
\begin{itemize}
  \item Ticketing-Daten als einzige Quelle (kein Join mit Monitoring im Release 1).
  \item Sechs der zehn Features.
  \item Baseline-Heuristik als Fallback bei Modell-Ausfall.
  \item Dashboard mit den sechs Kernkacheln.
\end{itemize}

\textbf{Bewusst nicht im MVP:}
\begin{itemize}
  \item Monitoring-Signal-Join (komplex, schwankende Datenqualität).
  \item Customer-Tier-Feature (Sensibilität gegenüber Fairness/Bias prüfen).
\end{itemize}

\textbf{Kompensation:} manuelle ``Service Impact''-Kennzeichnung im Ticket; wöchentliche Outcome-Reviews.

\subsection{Annahmen und Frühwarnsignale}

\begin{itemize}
  \item \textbf{Annahme:} Baseline-Heuristik liefert für Training/Start brauchbare Proxy-Labels.
  \item \textbf{Annahme:} Mit 6 Features ist Genauigkeit $>$75\,\% erreichbar.
  \item \textbf{Frühwarnsignal 1:} Alarmmüdigkeit (Acknowledge-Rate sinkt unter 60\,\%) $\rightarrow$ Schwellen neu justieren.
  \item \textbf{Frühwarnsignal 2:} False Negatives nehmen zu (SLA-Breaches ohne vorigen Alert) $\rightarrow$ Modell re-trainieren oder Features ergänzen.
\end{itemize}

\subsection{Lessons aus der Fallstudie}

\begin{itemize}
  \item Baseline-Heuristik ist ein vollwertiger MVP, nicht nur eine ``Notlösung''.
  \item Daten-Joins früh planen, aber nicht früh erzwingen.
  \item Anti-Gaming und Alarmmüdigkeit gehören in das Design, nicht in das Post-Mortem.
\end{itemize}

% --- 8. MANAGEMENT-PERSPEKTIVE -------------------------------------------
\section{Management-Perspektive: Analytics als verantwortetes Produkt}

Auf Wissens-Ebene sind Enhancement und Predictive methodische Themen. Auf Anwendungs-Ebene sind sie Handwerk. Auf Management-Ebene werden sie zu einem \emph{Analytics-Produkt} -- mit Product Owner, Releases, Monitoring und Audit. \quelle{Davenport, 2018; Shmueli, 2010}

\subsection{Zwei zentrale Zielkonflikte}

\begin{enumerate}
  \item \textbf{Performance vs.\ Erklärbarkeit.} Hochgenaue Black-Box-Modelle sind oft schwer zu erklären -- problematisch im regulierten Umfeld, in dem Entscheidungen begründet werden müssen. Praktische Lösung: \emph{erklärbares Modell mit etwas niedrigerer Genauigkeit für Produktionseinsatz, Black-Box-Modell als Vergleichs-Benchmark.}
  \item \textbf{Transparenz vs.\ Vertrauen.} Mehr Daten und Vorhersagen können als Überwachung empfunden werden. Praktikabel: aggregierte Sichten priorisieren, individuelle Vorhersagen mit eindeutigem Handlungs-Nutzen statt mit ``Score'' kommunizieren, Betriebsrat und Datenschutz früh einbinden.
\end{enumerate}

\subsection{Analytics als Produkt -- vier Pflichtelemente}

\begin{itemize}
  \item \fachbegriff{Product Owner.} Verantwortet das Analytics-Produkt fachlich und priorisiert die Pipeline.
  \item \fachbegriff{Release-Zyklus.} Modelle haben Lebenszyklus -- Training, Validierung, Shadow-Mode, Production, Retirement.
  \item \fachbegriff{Modell-Monitoring.} Vorhersage-Qualität wird selbst beobachtet (Genauigkeit, Drift, Bias). Auch das Modell hat ``Operations''.
  \item \fachbegriff{Audit-Trail.} Welche Vorhersage zu welcher Entscheidung geführt hat, ist nachvollziehbar -- in regulierten Umfeldern Pflicht.
\end{itemize}

\subsection{Konzeptdrift und Modell-Operations}

ML-Modelle altern. ``Konzeptdrift'' bedeutet, dass die Realität sich gegenüber den Trainingsdaten verschiebt, ohne dass es jemand bemerkt. Drei Warnzeichen:

\begin{itemize}
  \item Genauigkeit sinkt schleichend.
  \item Vorhersage-Verteilungen verschieben sich (mehr High-Risk-Cases, ohne dass mehr Breaches eintreten).
  \item Feature-Verteilungen ändern sich (neue Service-Kategorien, andere Schicht-Muster).
\end{itemize}

Senior-Praxis: Modell-Monitoring von Beginn an, mit klaren Re-Trainings-Triggern (z.\,B.\ ``Genauigkeit sinkt unter 70\,\% für zwei Wochen $\rightarrow$ Re-Training''.).

\subsection{Datenethik und Governance}

\begin{itemize}
  \item \fachbegriff{Fairness/Bias.} Wenn historische Daten Diskriminierung enthalten, lernt das Modell sie. Pflicht-Check auf protected attributes (z.\,B.\ Geschlecht, Alter, ethnische Zugehörigkeit), wo relevant.
  \item \fachbegriff{Erklärbarkeit.} Für Entscheidungen mit Außenwirkung (z.\,B.\ Bonität, Recruiting) ist Erklärbarkeit häufig regulatorisch gefordert.
  \item \fachbegriff{Zugriff und Audit-Trail.} Wer hat welche Vorhersage gesehen, mit welcher Begründung gehandelt? -- Pflicht in regulierten Umfeldern.
\end{itemize}

\subsection{``Minimum Viable Data Product''}

Senior-Disziplin: \emph{lieber ein kleines Data Product, das wirklich verwendet wird, als ein großes, das niemand benutzt.} Drei Kriterien für das MVP:

\begin{itemize}
  \item Eine klare, dokumentierte Entscheidung, die unterstützt wird.
  \item Ein nachvollziehbares Baseline-Modell oder eine erklärbare Heuristik.
  \item Ein Playbook, das jeden Alarm mit Handlung verknüpft.
\end{itemize}

\subsection{Entscheidungen unter Unsicherheit}

\begin{itemize}
  \item \textbf{Welche drei Use Cases zuerst -- trotz Datenlücken?} Schriftlich entscheiden, verworfene Alternative dokumentieren, Frühwarnsignal benennen.
  \item \textbf{Welche KPIs werden ``Executive''?} Bewusst auswählen, Anti-Gaming-Mechanismus festlegen, Review-Rhythmus.
\end{itemize}

\begin{merkbox}[Schluss-Merksatz für Tag 8]
Predictive ohne Playbook ist Reporting; Playbook ohne Owner ist Theater; Owner ohne Modell-Monitoring ist Hoffnung. Senior-Steuerung verbindet alle vier: \emph{Daten $\rightarrow$ Modell $\rightarrow$ Alert $\rightarrow$ Owner $\rightarrow$ Playbook $\rightarrow$ Modell-Monitoring.}
\end{merkbox}

% --- 9. TAKE-AWAYS -------------------------------------------------------
\section{Take-aways aus Tag 8}

\begin{enumerate}
  \item \textbf{Drei Reifegrade.} Reporting (Vergangenheit) $\rightarrow$ Real-time (Gegenwart) $\rightarrow$ Predictive (Zukunft).
  \item \textbf{Enhancement schließt den Lifecycle.} Discovery + Conformance + Enhancement zusammen.
  \item \textbf{Baseline vor ML.} Erklärbare Regeln bringen oft 60--80\,\% des Nutzens, sofort.
  \item \textbf{Use Case Canvas.} Ziel, Prognose, Handlung, Datenminimum, Risiko, Erfolg -- alle sechs Felder gefüllt.
  \item \textbf{O-H-E-D-P.} Owner, Handlung, Eskalation, Dokumentation, Post-Review pro Alert-Playbook.
  \item \textbf{ETL ist 60\,\% des Wertes.} Mapping, Case-ID-Strategie, Quality Checks -- die nicht-glanzvolle, aber tragende Arbeit.
  \item \textbf{Dashboards mit Handlung.} Trend + Ziel + Segmentierung + Next-Best-Action.
  \item \textbf{Anti-Gaming.} Jede KPI mit Gegenkontrolle, jede problematische KPI bewusst weggelassen.
  \item \textbf{Konzeptdrift überwachen.} Modell-Operations als eigene Disziplin -- ML altert.
\end{enumerate}

% --- 10. LITERATUR -------------------------------------------------------
\clearpage
\section{Literatur}

Im Skript verwendete und für die Vertiefung empfohlene Quellen. Zitierstil: APA-7.

\begin{description}[style=nextline,leftmargin=18pt,labelindent=0pt,itemsep=4pt]
  \item[Davenport, T.\,H. (2018).] \emph{The AI advantage: How to put the artificial intelligence revolution to work}. MIT Press.
  \item[Few, S. (2013).] \emph{Information dashboard design: Displaying data for at-a-glance monitoring} (2nd ed.). Analytics Press.
  \item[Inmon, W.\,H. (2005).] \emph{Building the data warehouse} (4th ed.). Wiley.
  \item[Kimball, R., \& Ross, M. (2013).] \emph{The data warehouse toolkit: The definitive guide to dimensional modeling} (3rd ed.). Wiley.
  \item[Maggi, F.\,M., Di Francescomarino, C., Dumas, M., \& Ghidini, C. (2014).] Predictive monitoring of business processes. In M.~Jarke et al.\ (Hrsg.), \emph{Advanced information systems engineering (CAiSE 2014)} (Lecture Notes in Computer Science, Vol.~8484, S.~457--472). Springer. \href{https://doi.org/10.1007/978-3-319-07881-6_31}{https://doi.org/10.1007/978-3-319-07881-6\_31}
  \item[Russom, P. (2011).] \emph{Big data analytics} (TDWI Best Practices Report, Q4 2011). The Data Warehousing Institute. \href{https://tdwi.org/research/2011/09/best-practices-report-q4-big-data-analytics.aspx}{https://tdwi.org/research/2011/09/best-practices-report-q4-big-data-analytics.aspx}
  \item[Shmueli, G. (2010).] To explain or to predict? \emph{Statistical Science, 25}(3), 289--310. \href{https://doi.org/10.1214/10-STS330}{https://doi.org/10.1214/10-STS330}
  \item[van der Aalst, W.\,M.\,P. (2016).] \emph{Process mining: Data science in action} (2nd ed.). Springer. \href{https://doi.org/10.1007/978-3-662-49851-4}{https://doi.org/10.1007/978-3-662-49851-4}
\end{description}

% --- 11. GLOSSAR ---------------------------------------------------------
\clearpage
\section{Glossar}

Die wichtigsten Begriffe des Tages -- kurz definiert, in alphabetischer Reihenfolge.

\begin{description}[style=nextline,leftmargin=0pt,labelindent=0pt,itemsep=4pt]
  \item[Alert] Auslöser einer Reaktion, wenn ein KPI oder ein Risiko-Score eine Schwelle überschreitet.
  \item[Alarmmüdigkeit] Effekt, bei dem zu viele Alerts dazu führen, dass Bearbeitende sie ignorieren.
  \item[Anomalieerkennung] ML-Disziplin zur Identifikation ungewöhnlicher Pfade, Daten oder Verhaltensmuster.
  \item[Baseline-Heuristik] Einfache, erklärbare Entscheidungsregel, oft als Vorstufe oder Vergleich zu ML.
  \item[BI-Dashboard] Visualisierungs-Werkzeug, das KPIs, Trends und Alerts für eine Zielgruppe konsolidiert.
  \item[Clustering] ML-Disziplin zur Gruppierung ähnlicher Cases ohne vorgegebene Labels.
  \item[Data Quality Check] Regel-Test, der Daten vor dem Laden ins DWH auf Vollständigkeit, Konsistenz und Plausibilität prüft.
  \item[DWH] \emph{Data Warehouse.} Konsolidierter Speicher analytischer Daten, häufig mit dimensionalem Modell.
  \item[Dimension] Beschreibende Attribute in dimensionaler Modellierung (Zeit, Kunde, Produkt, Service).
  \item[Drift (Konzeptdrift)] Schleichende Veränderung der Realität gegenüber den Trainingsdaten eines ML-Modells.
  \item[ETL/ELT] \emph{Extract-Transform-Load} bzw.\ \emph{Extract-Load-Transform.} Verfahren zur Befüllung eines DWH.
  \item[Enhancement] Disziplin im Process Mining: Anreicherung von Modellen mit Daten (Zeiten, Kosten, Ressourcen).
  \item[Faktentabelle] Zentrale Tabelle in dimensionaler Modellierung; jede Zeile ein Event mit Maßzahlen und Fremdschlüsseln.
  \item[Fairness/Bias] Eigenschaft eines ML-Modells, geschützte Attribute nicht zu Diskriminierung zu verwenden.
  \item[Klassifikation] ML-Disziplin: Vorhersage einer kategorialen Größe (ja/nein, A/B/C).
  \item[Konzeptdrift-Monitoring] Beobachtung der Modell-Genauigkeit und Feature-Verteilungen, mit Re-Training-Trigger.
  \item[Minimum Viable Data Product] Kleinste sinnvolle Daten- oder Analytics-Lösung, die in Produktion verwendet wird.
  \item[Playbook] Strukturierte Handlungsanweisung für einen bestimmten Alarm: Owner, Handlung, Eskalation, Dokumentation, Post-Review.
  \item[Predictive Process Monitoring] ML-basierte Vorhersagen auf Prozessdaten (SLA, Rework, Abbruch).
  \item[Real-time Monitoring] Beobachtung von Prozessen in Near-Real-Time mit Alerts und Dashboards.
  \item[Regression] ML-Disziplin: Vorhersage einer kontinuierlichen Größe (Dauer, Kosten).
  \item[Staging] Bereich, in dem Rohdaten unverändert vor der Transformation abgelegt werden.
  \item[Star-Schema] Dimensionales Datenmodell mit zentraler Faktentabelle und umlaufenden Dimensionen.
  \item[Use Case Canvas] Struktur-Werkzeug zur Definition eines Predictive Use Cases (Ziel, Prognose, Handlung, Daten, Risiko, Erfolg).
\end{description}

\vspace{1cm}
\noindent{\color{lpAccent}\rule{\linewidth}{0.6pt}}\\[4pt]
{\footnotesize\sffamily\color{lpGray}Lernplatt \textperiodcentered{} by Can Siebert \textperiodcentered{} Kurs 71 (Dig 42) \textperiodcentered{} Tag 8 \textperiodcentered{} \textcopyright{} 2026}

\end{document}
